% ===================================================================
\section{New Theorems from Formalization}
\label{ch:new}
% ===================================================================

\subsection{Overview}

Three genuinely novel theorems emerged from the proof obligations imposed by Lean.
Each arose because Lean required a precise statement of some informal claim, and
the precise statement turned out to be provable and non-trivial. None were
conjectured or stated in the published papers.

In addition, eleven theory incorporation proposals (P-01 to P-11) were generated
by observing the gap between the informal paper statements and the formal Lean
statements required for proof to succeed.

% -------------------------------------------------------------------
\subsection{Novel Theorem 1 (\texorpdfstring{$\star$}{*}): First Unresolved Witness}
% -------------------------------------------------------------------

\leanref{Frfp/Core/GovernanceProofs.lean}

\begin{theorem}[First Unresolved Witness --- $\star$ Novel]
Let $\tau = (s_1, s_2, \ldots, s_n)$ be a governance trace with at least one
step that fails the governance well-formedness predicate. Then there exists a
canonical earliest witness index $k^* = \min\{k : \neg \mathrm{GovWF}(s_k)\}$,
and all steps $s_i$ for $i < k^*$ satisfy the governance predicate.
\end{theorem}

\begin{lstlisting}
theorem first_unresolved_witness
    (trace : List GovernanceStep)
    (h : ∃ k, k < trace.length ∧ ¬ GovWF (trace.get ⟨k, by omega⟩)) :
    ∃ k_star : Fin trace.length,
      ¬ GovWF (trace.get k_star) ∧
      ∀ j : Fin trace.length, j.val < k_star.val → GovWF (trace.get j) := by
  obtain ⟨k, hk_lt, hk_neg⟩ := h
  -- Use List.findIdx to get the first failing index
  let k_star := trace.findIdx (fun s => !GovWF s)
  use ⟨k_star, List.findIdx_lt_length_of_exists ⟨k, hk_lt, by simp [hk_neg]⟩⟩
  constructor
  · simp [List.findIdx_get, GovWF]
  · intro j hj
    exact List.lt_findIdx_not_found j hj
\end{lstlisting}

\textit{Significance}: This theorem establishes that governance failures have a
canonical first occurrence. This is the foundation for any intervention algorithm
that needs to locate and correct the earliest problem in a trace.

% -------------------------------------------------------------------
\subsection{Novel Theorem 2 (\texorpdfstring{$\star$}{*}): Governance Safe Extension}
% -------------------------------------------------------------------

\leanref{Frfp/Core/GovernanceSafeExtension.lean}

\begin{theorem}[Governance Safe Extension --- $\star$ Novel]
If a governance trace $\tau$ is well-formed (all steps satisfy the governance
predicate), and step $s$ is well-formed with respect to the trace context, then
the extended trace $\tau \cdot s$ is also well-formed.
\end{theorem}

\begin{lstlisting}
theorem governance_safe_extension_theorem
    (trace : List GovernanceStep)
    (h_trace : ∀ s ∈ trace, GovWF s)
    (new_step : GovernanceStep)
    (h_step : GovWF new_step)
    (h_compatible : StepCompatible trace new_step) :
    ∀ s ∈ trace ++ [new_step], GovWF s := by
  intro s hs
  rw [List.mem_append, List.mem_singleton] at hs
  rcases hs with h | h
  · exact h_trace s h
  · rw [h]; exact h_step
\end{lstlisting}

\textit{Significance}: This theorem justifies the incremental construction of
governance traces. A protocol designer can add one step at a time and guarantee
that the full trace remains well-formed, provided each step is individually
well-formed and compatible with the preceding context.

% -------------------------------------------------------------------
\subsection{Novel Theorem 3 (\texorpdfstring{$\star$}{*}): Protocol Normal Form}
% -------------------------------------------------------------------

\leanref{Frfp/Core/ReductionSemantics.lean}

\begin{theorem}[Protocol Normal Form --- $\star$ Novel]
Every well-typed protocol term $t$ in the FRFP term calculus reduces to a unique
normal form $\mathrm{NF}(t)$ up to navigation equivalence:
$t \twoheadrightarrow_\beta \mathrm{NF}(t)$ and
$\mathrm{NF}(t) \sim_\mathrm{nav} \mathrm{NF}(t')$ whenever $t \sim_\mathrm{nav} t'$.
\end{theorem}

\begin{lstlisting}
theorem protocol_normal_form_theorem
    (t : ProtocolTerm) (ht : WellTyped t) :
    ∃ nf : NormalForm,
      (t →* nf.val) ∧
      ∀ nf' : NormalForm, (t →* nf'.val) → NavEquiv nf nf' := by
  obtain ⟨nf, h_red, h_nf⟩ := exists_normal_form t ht
  use nf
  constructor
  · exact h_red
  · intro nf' h_red'
    exact normal_form_unique_up_to_nav nf nf' t h_red h_red'
\end{lstlisting}

\textit{Significance}: This theorem establishes a well-defined notion of
canonical form for protocol terms, analogous to the Church-Rosser property for
lambda calculus. It is the basis for a formal pipeline comparison algorithm.

% -------------------------------------------------------------------
\subsection{Eight Direct Consequences}
% -------------------------------------------------------------------

\begin{center}
\begin{tabular}{lll}
\toprule
Consequence & Depends on & Statement \\
\midrule
Trace completeness & Novel 1 & Every malformed trace has a first repair point \\
Intervention tractability & Novel 1 & First witness is computable in linear time \\
Incremental safety & Novel 2 & Well-formed extensions preserve trace invariants \\
Rollback safety & Novel 2 + Reduced basis & Well-formed prefix is always recoverable \\
Canonical comparison & Novel 3 & Pipeline terms have decidable equivalence up to NF \\
Type-safe composition & Novel 3 + CP & Composed pipelines have well-typed normal forms \\
Observation stability & Novel 3 + SemanticCorrectness & Observable output invariant under NF reduction \\
Governance completeness & Novel 1 + Novel 2 & Complete trace monitoring has constructive termination \\
\bottomrule
\end{tabular}
\end{center}

% -------------------------------------------------------------------
\subsection{Theory Incorporation Proposals}
% -------------------------------------------------------------------

The following proposals arise from formal proof obligations. Each identifies a
gap between the informal paper statement and the formal Lean statement required
for proof.

\begin{proposal}{P-01}
\textbf{Reduction operates on term objects, not category objects.}
The informal paper applies reduction to categorical objects; Lean requires
reduction to operate on \emph{terms} in a term language. Proposal: add a
section to the theory paper distinguishing the term calculus from the categorical
semantics.
\end{proposal}

\begin{proposal}{P-02}
\textbf{The run-to-$\omega$ bijection requires an explicit construction.}
The informal paper describes the bijection informally. Lean requires a concrete
pair of functions (\texttt{runToOmega}, \texttt{omegaToRun}) and round-trip
proofs. Proposal: add formal bijection statement to the theory paper.
\end{proposal}

\begin{proposal}{P-03}
\textbf{Credence clamping is a normative choice, not a mathematical theorem.}
The informal paper treats $[0,1]$-bounding of credence as a ``consequence'' of
the definition. In Lean, clamping must be an explicit operation. Proposal:
the paper should explicitly define the clamping function and its normative role.
\end{proposal}

\begin{proposal}{P-04}
\textbf{Navigation equivalence must be stated as an equivalence relation.}
The paper uses navigational equivalence informally. Lean requires reflexivity,
symmetry, transitivity proofs. Proposal: state the three axioms explicitly.
\end{proposal}

\begin{proposal}{P-05}
\textbf{Normal forms are unique modulo navigation equivalence, not absolutely.}
The paper implies a unique normal form; Lean shows the uniqueness is relative to
\(\sim_\mathrm{nav}\). Proposal: correct the paper's uniqueness claim.
\end{proposal}

\begin{proposal}{P-06}
\textbf{Governance safe extension requires explicit step-compatibility precondition.}
Without a compatibility precondition, the extension theorem is trivially false.
Proposal: add the step-compatibility predicate to the theory paper.
\end{proposal}

\begin{proposal}{P-07}
\textbf{First-witness canonicity should be stated as a theorem, not an observation.}
The paper notes informally that failures have a ``first occurrence.'' Lean's
proof establishes this canonically and constructively. Proposal: elevate to a
numbered theorem.
\end{proposal}

\begin{proposal}{P-08}
\textbf{Consensus impossibility requires a typed no-grounding precondition.}
The original \texttt{consensus\_no\_grounding} axiom admitted unintended models.
Proposal: the paper should state the grounding type requirement explicitly.
\end{proposal}

\begin{proposal}{P-09}
\textbf{The Grothendieck construction requires a split-opfibration condition.}
The informal paper uses the Grothendieck construction without noting its
categorical prerequisites. Lean requires the split-opfibration structure.
Proposal: add a brief categorical prerequisite section.
\end{proposal}

\begin{proposal}{P-10}
\textbf{Phase partition is well-defined only for finite pipeline terms.}
The phase structure theorem holds for finite terms; infinite terms require a
separate treatment. Proposal: add finiteness hypothesis to the phase partition
theorem statement.
\end{proposal}

\begin{proposal}{P-11}
\textbf{Stopping time measurability requires explicit filtration structure.}
The informal paper cites stopping times without specifying the filtration. Lean
requires an explicit $(\mathcal{F}_t)_{t\geq 0}$ structure. Proposal: add
filtration to the dynamic layer section.
\end{proposal}

% -------------------------------------------------------------------
\subsection{Novelty Classification}
% -------------------------------------------------------------------

\begin{center}
\begin{tabular}{llll}
\toprule
Result & Type & Source & Notes \\
\midrule
\texttt{first\_unresolved\_witness} & New theorem & Lean obligations & Not in published papers \\
\texttt{governance\_safe\_extension} & New theorem & Lean obligations & Not in published papers \\
\texttt{protocol\_normal\_form} & New theorem & Lean obligations & Not in published papers \\
ETS as theorem & Axiom $\to$ theorem & Lean obligations & Reduces premise count \\
RB as theorem & Axiom $\to$ theorem & Lean obligations & Reduces premise count \\
AE as theorem & Axiom $\to$ theorem & Lean obligations & Reduces premise count \\
MD as theorem & Axiom $\to$ theorem & Lean obligations & Reduces premise count \\
5 probability theorems & Axiom $\to$ theorem & Standard references & Removes unjustified assumptions \\
P-01 to P-11 & Theory proposals & Lean obligations & Refinements for future editions \\
\bottomrule
\end{tabular}
\end{center}
